\begin{exercisechunk}
\item \textbf{Continued training on incomplete data.}
\label[exercise]{ex:l20-missing-contexts}
\exerciseprofile{2}{2}{2}{\exproof\quad\excalculation\quad\exinterpretation}{%
Distinguish drift on unobserved contexts from a zero sampled update on successful data.}
Use the binary logistic policy and deployment population of
\cref{ex:missing-contexts}: $x_1=(1,0)^\top$ and $x_2=(-1,1)^\top$
are equally likely and both have label $1$. Training uses only $x_1$.
Write $\theta=(u,v)^\top$ and start at $(M,2M)^\top$ for $M>0$.
\begin{enumerate}[(a)]
\item Derive the flows for $L_{\rm tr}=-\log\sigma(u)$ and
$J_{\rm tr}=\sigma(u)$, and show that in both cases
$v(t)=2M$ and $u(t)\to\infty$.
Deduce that the initial deterministic classifier is correct on both
contexts, but eventually misclassifies $x_2$.
Show that deployment reward starts at $\sigma(M)$ and tends to $1/2$,
although the deployment population is linearly separable.
\item At a fixed parameter, let $A_1,\ldots,A_n$ be independent draws
from a fixed distribution $q$ with $q(-1),q(1)>0$.
Let $R_i=\I\{A_i=1\}$ and define
\[
 \widehat g=\frac1n\sum_{i=1}^n
 \frac{\pi_\theta(A_i\mid x_1)}{q(A_i)}
 (R_i-1)\nabla\log\pi_\theta(A_i\mid x_1).
\]
Prove that $\E[\widehat g]=\nabla J_{\rm tr}(\theta)$, although
$\widehat g=0$ whenever all $R_i=1$.
Explain why the expectation identity need not hold after conditioning
on all rewards being one, and why multiplying the retained terms by
finite weights cannot restore it.
\end{enumerate}
\end{exercisechunk}
